\documentclass[10pt,a4paper]{amsart}
\usepackage[margin=19mm]{geometry}
\usepackage{amsmath,amssymb,mathtools}
\usepackage[scr=rsfs,cal=euler]{mathalfa}
\usepackage{xfrac}
\usepackage[unicode,hidelinks,bookmarks=false]{hyperref}
\newcommand{\ie}{i.e.\ }
\newcommand{\cf}{cf.\ }
\newcommand{\resp}{respectively\ }
\newcommand{\Subsection}{\subsection}
\providecommand{\nobreakdash}{\nobreak\hbox{-}}
\setlength{\emergencystretch}{2em}
\multlinegap0pt
\usepackage{float}
\usepackage{amssymb}
\usepackage{xcolor}
\newtheorem{lemma}{Lemma}
\title{Uniqueness of Flotation and Buoyancy Surfaces for Convex Polytopes}
\author{S. Dann, O. Herscovici, S. Myroshnychenko}
\date{Source: arXiv:2605.09776v1 (CC BY 4.0)}
\begin{document}
\maketitle

\noindent\emph{Source and licence.} Uniqueness of Flotation and Buoyancy Surfaces for Convex Polytopes. Version arXiv:2605.09776v1 (CC BY 4.0), distributed by arXiv under the Creative Commons Attribution 4.0 International licence, http://creativecommons.org/licenses/by/4.0/, as recorded in the arXiv licence metadata for this exact version and shown on the abstract page https://arxiv.org/abs/2605.09776v1. Full licence notice: \texttt{licenses/arxiv-2605.09776-CC-BY-4.0.txt}.\par\smallskip

\noindent\textit{Editorial scope.} Counted unit: the article's lemma lemma (source0.tex lines 350--352). The article's complete original proof of it is delivered with it (source0.tex lines 353--395, one \texttt{proof} environment of 43 lines). No mathematics is written, repaired or re-derived: the statement and the proof are the article's own lines, in the article's own notation, and no retained line is substituted or re-flowed.  Retained uncounted as the source-local context the proof needs: lines 317--323; lines 334--336.  One declared adaptation: exactly 1 ancillary strings inside the retained spans are omitted (image-inclusion commands and, where the source repeats a label, the repeated label definitions) and no other line differs from the source; the figure environments, with their placement, captions and labels, are kept, and the package records the omitted strings in fidelity receipts bound to the affected bodies.  No retained line contains a citation, so the wrapper carries no bibliography and no bibliography entry.  The retained lines contain the article's own diagram code, which the wrapper typesets with the standard tikz packages; no image file is delivered and the package declares no asset dependency.  Editorial formatting only, in addition: an AMS \texttt{amsart} preamble that restates the article's own declarations and macros used by the retained lines, namely the environments \texttt{lemma} on separate counters, together with the standard packages those lines need.  The retained environments print in this document as Lemma 1 in order of appearance, whereas the article prints the same items with the numbering of its own full document.  The complete native source of the article as distributed in the arXiv e-print for this exact version is bundled unchanged as \texttt{sources/200256/source0.tex}.
\begin{lemma}\label{P_Dupin}
For a polygon $P$, $P_{[\delta]}$ is a finite union of hyperbolic arcs. At a common point $p$ of two consecutive hyperbolic arcs we have one of the following:
\begin{description}
    \item[Case 1] Hyperbolic arcs have a common tangent line (smooth union). This tangent line to $P_{[\delta]}$ at $p$ intersects two non-parallel sides of~$P$.
    \item[Case 2] Hyperbolic arcs have distinct tangent lines (non-smooth union). The set of supporting lines to $P_{[\delta]}$ at $p$ intersects a pair of parallel sides of $P$. 
\end{description}
\end{lemma}

\begin{lemma}\label{Sus}
    In $\mathbb{R}^d$, $d \geq 2$, $\delta \neq \frac{1}{2}$ if and only if for any point $x \in \partial K$ and $(d-2)$-dimensional subspace $\ell$, $x \in \ell$, $\ell \cap \text{int}K = \emptyset$, there exist two distinct hyperplanes that support $K_{[\delta]}$ and contain $\ell$.
\end{lemma} 

\begin{lemma}
    For $\delta \neq \frac{1}{2}$ and a given $P_{[\delta]}$, there exists a unique polygon $P$ that can be reconstructed from~$P_{[\delta]}$.
\end{lemma}

\begin{proof}
    To reconstruct $P$, we first need to understand how the supporting lines of $P_{[\delta]}$ intersect sides of $P$. As Lemma \ref{P_Dupin} explains, $P_{[\delta]}$ is a finite union of hyperbolic arcs that are joined at the end points in a smooth or non-smooth way. We start by determining all asymptotes of the hyperbolic arcs of $P_{[\delta]}$, and considering the intersection $W$ of all quarter-planes containing the corresponding arcs of hyperbolas. Note that $W$ is also a polygon. By the above observation, $P \subseteq W$ and $\partial P \cap \partial W \neq \emptyset$, where the intersection consists of a finite number of segments. Note that by convexity, each subset of $\partial P \cap \partial W$ that belongs to one line can be assumed to be a segment.
    
    To determine missing sides of $P$, we plot the one-sided tangent rays to non-smooth adjoint hyperbolic arcs at the common points. The sides of $W$ intersected by these two rays must contain the end points of missing sides; the segments between the intersected points on such sides of $W$ are subsets of $\partial P$.
    
    Observe that the intersection of a supporting line of $P_{[\delta]}$ with $\partial P$ is continuous with respect to a direction vector $\theta \in \mathbb{S}^2$.
    We partition each side of $P$ into a finite number of open segments: we exclude all vertices and all points on the sides that lie on liquid levels passing through vertices.
    The number of such excluded points is bounded by three times the number of vertices of $P$ (since each vertex may exclude at most two distinct points).
    This divides $\partial P$ into a finite number of segments. 
    Thus, for any point $x$ in a given segment $(pq)$, two supporting lines of $P_{[\delta]}$ that contain $x$ intersect the same pair of segments (which may coincide) regardless of the choice of $x \in (pq)$.

    
    For $x \in (pq)$, by Lemma~\ref{Sus}, let $y \in (rs)$ and $z \in (uv)$ be the two distinct points where the lines of liquid through $x$ intersect $\partial P$. 
    Then we investigate the following cases based on the mutual position of the intersected segments $(pq), (rs), (uv)$.

       \begin{figure}[h!]
        \centering
        
         \caption{For $t \in P_{[\delta]}$, case $(rs)\neq (uv)$ and $pq \parallel (uv)$.}
    \label{fig:2DParallel}
     \end{figure}
     
    \begin{itemize}
     \item If $(rs)\neq (uv)$ and $(pq) \not \parallel (uv)$, then $(pq), (uv) \subset \partial W$.

     \item If $(rs)\neq (uv)$ and $(pq) \parallel (uv)$, then $(pq)$ may be parallel to $(rs)$, Figure~\ref{fig:2DParallel}. Then, $(rs)$ and $(uv)$ belong to the same side of $P$ (otherwise, they are not parallel and belong to $\partial W$). 
     If we let $x$ approach $q$ (then $y$ approaches $\tilde{y} \in [rs)$), then the second line of liquid through $y$ may intersect the same side to which $(pq)$ belongs.
     
     The intersection cannot be a point $w_1$ to the left from $x$ along the line through $p, q$, since then the area to the left from the line through $w_1, y$ is strictly less than $\delta$ (less than the area on the left from line through $x, y$). 
     If the intersection is at a point $w_2 \in (pq]$, then when $x$ approaches $q$, $w_2$ also approaches $q$, so $x=w_2=q$ in the limit. However, Lemma \ref{Sus} guarantees $x\neq w_2$. 
     Thus, the point is in a position of $w_3$ to the right from $q$ along the line through $p,q$.
    Consequently, $w_3$ lands in a segment $(p_1q_1) \neq (pq)$ on the right from $(pq)$. Hence, we can repeat the same considerations starting from $(p_1q_1)$ until we exhausted all segments laying on the same side as $(pq)$ or $(rs)$. 
    As the number of such segments on each side is finite, the point corresponding to $\tilde{y}$ or $w_3$ will belong to a segment not parallel to $(pq)$.
    Hence, that segment and its non-parallel predecessor $(s_1s_2)$ belong to $\partial W$. Based on $P_{[\delta]}$ and a given segment $(s_1s_2) \in \partial P$, we reconstruct the previous segments all the way to $(pq)$.
    

     \item If $(rs)=(uv)$ and $(pq) \not \parallel (uv)$, analogously $(pq), (uv) \subset \partial W$

     \item If $(rs)= (uv)$, and $(pq) \parallel (uv)$, we can apply the same approach as $(rs) \neq (uv)$ and $(pq) \parallel (uv)$, because the position of $z \in (rs)$ does not affect the argument for $x$ approaching $q$.
    \end{itemize}

We conclude that by continuously varying the liquid line to intersect $\partial W$, we reconstruct $\partial P$, except for a finite number of points that we excluded. Taking the closure of this set yields $\partial P$, and $P$.
\end{proof}

\end{document}
